\chapter{What Makes a Language Model Large}
\label{ch:what-makes-language-models-large}

\abstract{Large language models are not defined by parameter count alone. They are coupled artifacts whose behavior emerges from data selection, tokenization, architecture, optimization, post-training, inference systems, and evaluation. This chapter establishes that lifecycle view. It explains why causal prediction became the dominant training interface, why scaling laws changed the engineering economics of language modeling, why instruction following and preference learning turned pretrained models into assistant systems, and why reasoning models made inference-time computation part of the modeling problem.}

\section{The Object of Study}
\label{sec:object-of-study}

A language model assigns probabilities to token sequences, but a modern \llm is better understood as a trained conditional computation system whose distribution is shaped by corpus design, optimizer dynamics, architectural constraints, and deployment-time decoding. For a tokenized sequence $x_{1:T}$, the causal language-model factorization is
\begin{equation}
  p_{\theta}(x_{1:T}) = \prod_{t=1}^{T} p_{\theta}(x_t \mid x_{<t}),
\end{equation}
and the usual training objective minimizes the average negative log-likelihood,
\begin{equation}
  \mathcal{L}(\theta) = -\frac{1}{T}\sum_{t=1}^{T}\log p_{\theta}(x_t \mid x_{<t}),
\end{equation}
with perplexity defined as $\exp(\mathcal{L})$ when the loss is measured in natural-log units. The \transformer made this system practical at scale because self-attention exposes long-range token interactions while remaining highly parallel on accelerators \cite{vaswani2017attention}. GPT-style causal models then showed that a single left-to-right objective can support broad transfer when the model is trained on enough heterogeneous text \cite{radford2019language,brown2020language}. BERT-style masked models remain important for representation learning and retrieval, and dense retrieval systems later made representation quality central to open-domain question answering and RAG pipelines \cite{devlin2018bert,karpukhin2020dense,lewis2020retrieval}.

The word ``large'' therefore refers to several coupled axes, summarized in Table~\ref{tab:axes-of-largeness}. Parameter count changes capacity, but active parameters matter separately in sparse mixture-of-experts systems; token count changes the statistical evidence seen during training, but token provenance determines whether the evidence is useful or publishable; context length changes what can be conditioned on, but KV-cache memory and attention cost determine whether that context can be served economically \cite{kaplan2020scaling,hoffmann2022training,deepseek2024v3}. Inference budget is now also part of model behavior, because retrieval, reranking, sampling, verification, tool calls, and explicit reasoning traces can spend additional computation after pretraining has ended \cite{borgeaud2022retro,lewis2020retrieval,chen2025reasoning,yang2025testtime}.

\begin{table}[t]
\caption{Axes of largeness in modern language-model systems.}
\label{tab:axes-of-largeness}
\small
\begin{tabular}{@{}p{1.9cm}p{1.8cm}p{3.7cm}p{4.2cm}@{}}
\hline\noalign{\smallskip}
Axis & Symbol & What Scales & Typical Failure Mode \\
\noalign{\smallskip}\svhline\noalign{\smallskip}
Parameters & $N$ & Dense weights or total MoE weights & Capacity improves while memory, bandwidth, and overfitting risks grow \\
Active compute & $N_{\mathrm{active}}$ & Parameters used per token in sparse models & Routing imbalance or expert underuse hides nominal scale \\
Training data & $D$ & Tokens, documents, domains, and languages & More tokens amplify duplication, leakage, and licensing problems \\
Training compute & $C$ & FLOPs and accelerator time & Undertraining or wasted compute from poor data/model balance \\
Context & $T$ & Prompt and generated sequence length & KV-cache memory and attention cost dominate serving \\
Inference budget & $B$ & Samples, search, retrieval, tools, verifiers & Higher pass rates can hide weak calibration or invalid verifiers \\
\noalign{\smallskip}\hline
\end{tabular}
\end{table}

This book treats an \llm as a stack rather than as a single neural network. The stack begins with raw data and a tokenizer; passes through architecture, loss, optimizer, and distributed training; continues through supervised instruction tuning and preference optimization; and ends in an inference service that must decide how much context, retrieval, search, and safety filtering to apply. A mistake in any layer can be misread as a model limitation: a weak tokenizer can harm multilingual coverage, a contaminated benchmark can exaggerate capability, an unstable optimizer can waste compute, and a poorly designed serving path can make a strong model unusable under latency constraints \cite{shi2023detecting,xu2024benchmark}.

\begin{figure}[t]
\centering
\begin{tikzpicture}[node distance=7mm and 5mm]
\node[llmblock] (data) {Data\\provenance};
\node[llmblock, right=of data] (tok) {Tokenizer\\template};
\node[llmblock, right=of tok] (arch) {Architecture\\loss};
\node[llmblock, right=of arch] (train) {Optimization\\systems};
\node[llmblock, below=of train] (post) {Post-training\\alignment};
\node[llmblock, left=of post] (serve) {Inference\\service};
\node[llmblock, left=of serve] (eval) {Evaluation\\governance};
\draw[llmarrow] (data) -- (tok);
\draw[llmarrow] (tok) -- (arch);
\draw[llmarrow] (arch) -- (train);
\draw[llmarrow] (train) -- (post);
\draw[llmarrow] (post) -- (serve);
\draw[llmarrow] (serve) -- (eval);
\draw[llmdash] (eval.north) .. controls +(0,1.0) and +(0,-1.0) .. node[left, font=\scriptsize, align=center] {audit\\feedback} (data.south);
\node[llmnote, below=9mm of serve, text width=0.74\textwidth] (note) {A model release is a chain of contracts, not a checkpoint alone. Each layer inherits upstream assumptions and adds new failure modes: licensing and contamination in data, token efficiency in the tokenizer, numerical stability in training, behavioral shifts in post-training, and latency or policy failures in deployment.};
\end{tikzpicture}
\caption{The lifecycle stack used throughout the book. Each arrow represents a contract: the next layer inherits the assumptions, artifacts, and failure modes of the previous layer. The dashed loop emphasizes that evaluation and deployment incidents should feed back into data, training, and release practice rather than remain a final report.}
\label{fig:llm-lifecycle-stack}
\end{figure}

\section{Why Prediction Became a General Interface}
\label{sec:prediction-interface}

The central technical bargain of language modeling is that next-token prediction converts unstructured text into a dense self-supervised training signal.\index{next-token prediction}\index{self-supervised learning} Every token in a document can contribute a target, so web-scale corpora become usable without manual labels \cite{brown2020language}. The bargain is imperfect: the objective rewards distributional continuation, not truth, helpfulness, calibration, or harmlessness \cite{ouyang2022training}. Yet it is powerful because language co-occurs with facts, code, dialogue, procedures, arguments, and instructions, so a model trained to predict text can internalize many latent tasks before any task-specific fine-tuning \cite{brown2020language}.

The \transformer architecture amplified this bargain. Self-attention gives each position a data-dependent path to previous positions, so the model can condition on both local syntax and distant context \cite{vaswani2017attention}. Multi-head attention gives the model multiple learned projection subspaces for comparing and aggregating tokens, while feed-forward sublayers supply high-capacity tokenwise transformations. Residual pathways and normalization make very deep stacks trainable, and modern decoder families refine these ingredients with rotary position embeddings, RMS normalization, gated feed-forward networks, grouped-query attention, and mixture-of-experts routing \cite{su2021roformer,touvron2023llama,ainslie2023gqa,deepseek2024v3}.

Prediction also gives a clean implementation path. The model receives token ids, maps them to embeddings, repeatedly mixes information through attention and feed-forward blocks, and projects the final hidden states to vocabulary logits. Training minimizes cross-entropy between predicted logits and the next token, while generation samples or selects tokens from the resulting distribution. Later chapters derive this computation from tensor shapes and turn the derivations into original implementation exercises.

\section{Scaling Changed the Research Program}
\label{sec:scaling-changed-program}

Before the modern scaling era, natural language processing was organized around task-specific datasets, architectures, and losses. Scaling-law work reframed the field by showing that loss follows empirical power laws over model size, dataset size, and compute across broad ranges \cite{kaplan2020scaling}. Compute-optimal training then shifted attention from simply increasing parameters to balancing model size and token count, arguing that many early large models were undertrained relative to their compute budgets \cite{hoffmann2022training}. The consequence is practical: a serious training plan must budget parameters, tokens, sequence length, hardware throughput, and data quality together.

Open model reports made this systems view explicit, although they should be read as technical reports rather than as a substitute for independent replication. LLaMA demonstrated that carefully trained smaller models can be competitive when token budgets and data mixtures are chosen well \cite{touvron2023llama}. Llama 2 expanded the public discussion of pretraining, supervised fine-tuning, preference data, and safety evaluation \cite{touvron2023llama2}. Llama 3 pushed the recipe toward longer context, multilingual coverage, coding, safety models, and extensive post-training \cite{dubey2024llama3}. DeepSeek-V3 and Kimi K2 indicate that frontier-scale open-weight systems increasingly combine mixture-of-experts routing, specialized attention mechanisms, and agentic or coding-oriented post-training, but their claims still require careful benchmark and deployment-context interpretation \cite{deepseek2024v3,kimi2025k2}.

Scaling also exposed failure modes. Training loss can improve while benchmark validity degrades through leakage, contamination, or overfitting to public leaderboards \cite{shi2023detecting,xu2024benchmark}. Larger context windows can increase cost faster than they increase robust long-context reasoning. More samples at inference time can improve pass rates while hiding weak calibration or brittle verification. This book therefore treats evaluation as a modeling component, not as an afterthought.

\section{From Base Models to Assistants}
\label{sec:base-to-assistants}

A pretrained base model is not yet an assistant. It has learned broad continuation behavior, but user-facing behavior requires additional training signals that specify what counts as a helpful response. Supervised instruction tuning supplies demonstrations of desired behavior, and \rlhf adds preference rankings so the model can be optimized toward responses humans prefer under a specific annotation protocol \cite{ouyang2022training}. Direct preference optimization and related methods simplify this stage by fitting preferences without an explicit online reinforcement-learning loop, although they still depend on the quality and scope of the preference data \cite{rafailov2023direct}.

Parameter-efficient adaptation changed who can perform this second stage. LoRA freezes the base model and learns low-rank updates inside selected linear layers, which makes adaptation far cheaper than full fine-tuning for many settings \cite{hu2021lora}. QLoRA combines low-rank adaptation with quantized base weights, making instruction tuning feasible on much smaller hardware budgets \cite{dettmers2023qlora}. These methods are not magic compression; they change the trainable subspace, memory footprint, optimizer state, and failure modes. Chapter~\ref{ch:parameter-efficient-adaptation} treats rank, target modules, quantization error, and evaluation as design decisions rather than as fixed defaults.

Assistants also need retrieval and tools when the answer depends on fresh, private, or long-tail knowledge. Retrieval-augmented generation inserts external evidence into the context, while tool-use patterns let the model call search, calculators, databases, or code execution environments \cite{lewis2020retrieval,yao2023react}. These systems trade parametric memory for runtime grounding, but they introduce their own failure modes: retriever recall, reranker precision, context ordering, citation faithfulness, and prompt-injection risk.

\section{Reasoning Makes Inference Part of Training}
\label{sec:reasoning-inference-training}

Reasoning in \llms is not a single mechanism. Chain-of-thought prompting elicits intermediate natural-language steps, self-consistency samples multiple paths, tree search branches over candidate partial solutions, and verifier-guided methods score or filter reasoning traces \cite{wei2022chain,yao2023tree}. These traces can improve task performance without being faithful causal explanations of the model's internal computation, so they must be evaluated as generated artifacts rather than accepted as transparent reasoning records \cite{lyu2023faithful}. Reasoning-oriented reinforcement learning adds another layer by rewarding correct final answers, verified intermediate steps, or task-specific executable outcomes \cite{deepseek2025r1,wen2025rlvr}. A recurring post-2024 trend is that inference-time compute and training-time incentives must be studied together: some gains come from better policies, while others come from searching, sampling, verifying, or allocating more computation at test time \cite{chen2025reasoning,snell2024scaling,yang2025testtime}.

This perspective changes how we teach algorithms such as PPO, DPO, GRPO-style training, process reward models, and Monte Carlo tree search. The book does not present reinforcement learning as a detached prerequisite followed by an unrelated \rlhf recipe. Instead, it follows the sequence-level control problem: what is being rewarded, what distribution is being constrained, what verifier is trusted, and how much extra computation is spent at inference time.

\section{A Roadmap for the Book}
\label{sec:roadmap}

Part I defines the foundation: what is being modeled, how text becomes tokens, how corpora become training signals, and why provenance matters. Part II builds the architecture and systems spine, moving from attention and GPT-style decoders to LLaMA-class blocks, optimization, distributed training, and inference serving. Part III studies adaptation through instruction tuning, synthetic data, LoRA, QLoRA, domain adaptation, and multilingual specialization. Part IV covers retrieval, tools, agents, preference learning, reasoning, multimodality, evaluation, safety, and governance.

The intended reading discipline is simple. Each chapter starts from a concrete problem, states the mathematical object being optimized, connects the object to implementation, and then asks how the claim would be evaluated. That discipline is the main defense against treating slides, benchmark tables, or model cards as self-evident truth. The goal is not to memorize the names of model families; it is to understand why a modeling choice works, when it fails, and how to test it in a reproducible system.

\section{A Publication-Grade Reading Protocol}
\label{sec:publication-grade-reading-protocol}

A high-quality model book should teach readers how to interrogate claims. For any new model report, read the artifact as a tuple
\begin{equation}
  \mathcal{R} =
  (D, \tau, A, O, S, P, I, E, G),
\end{equation}
where $D$ is the data mixture, $\tau$ the tokenizer and templates, $A$ the architecture, $O$ the optimization recipe, $S$ the training system, $P$ the post-training process, $I$ the inference system, $E$ the evaluation evidence, and $G$ the governance constraints. A claim about capability is under-specified if any component is missing. For example, a long-context score without $I$ says little about KV-cache cost or scheduler feasibility; a reasoning score without $P$ and $I$ hides whether the gain came from reinforcement learning, longer traces, more samples, or a stronger verifier.

This protocol also explains why figures in this book are original synthesis diagrams rather than copied architecture art from papers. The purpose of a textbook figure is not to reproduce a model card screenshot; it is to expose the invariant structure that lets readers compare GPT decoders, MoE systems, RAG pipelines, multimodal connectors, and reasoning-time controllers under a single analytic vocabulary. When a diagram is inspired by a paper or system report, the caption cites the source; the drawing itself is a new explanatory abstraction.

\section{Key Terms}
\label{sec:chapter-one-key-terms}

\begin{description}[Inference budget]
\item[Base model] A pretrained model before instruction tuning, preference optimization, or task-specific adaptation.
\item[Context length] The number of tokens available to condition the model during a forward pass.
\item[Inference budget] The computation spent after a prompt is received, including decoding, retrieval, search, verification, and tool use.
\item[Perplexity] The exponential of average negative log-likelihood under a language model, useful for comparing models only when tokenization and evaluation data are controlled.
\item[Post-training] Training after base pretraining, including SFT, preference optimization, safety tuning, and domain adaptation.
\end{description}

\section{Exercises}
\label{sec:chapter-one-exercises}

\begin{enumerate}
\item Write the causal factorization and cross-entropy loss for a batch with shape $B \times T$. State which tokens should be masked for a chat-style SFT example.
\item Pick two rows from Table~\ref{tab:axes-of-largeness}. For each row, describe one benchmark result that could improve while the real system becomes worse.
\item Explain why a model with fewer total parameters can be more expensive to serve than a larger model under a different context length or KV-cache layout.
\item Design a provenance checklist for one dataset, one code repository, one figure, and one model checkpoint before they can be used in this book.
\item Give an example where a chain-of-thought answer is correct but the explanation is not a faithful record of the underlying computation.
\end{enumerate}
